\begin{table}[!htbp]
\centering
\caption{SAE low-income priority and the income gradient in school VA}
\label{tab:expected-va-priority}
\resizebox{\textwidth}{!}{%
\begin{tabular}{lccccccc}
\toprule
 & Math & Verbal & Exam taking ⊥ math & HP institution ⊥ math & HP field ⊥ math & Proj. income ⊥ math & Benefits app. ⊥ math \\
\midrule
\multicolumn{8}{l}{\textit{Panel A: First-choice VA}} \\
Income decile, no priority control & 0.094 & 0.111 & 0.092 & -0.021 & -0.011 & 0.011 & 0.069 \\
Income decile, with priority & 0.080 & 0.094 & 0.077 & -0.018 & -0.008 & 0.010 & 0.058 \\
SAE priority (low-income) & -0.110 & -0.129 & -0.115 & 0.028 & 0.025 & -0.010 & -0.089 \\
 & (0.005) & (0.004) & (0.004) & (0.003) & (0.004) & (0.004) & (0.004) \\
\midrule
\multicolumn{8}{l}{\textit{Panel B: Expected VA ($E_i^m$)}} \\
Income decile, no priority control & 0.064 & 0.087 & 0.088 & -0.014 & -0.015 & 0.007 & 0.056 \\
Income decile, with priority & 0.056 & 0.075 & 0.075 & -0.012 & -0.012 & 0.006 & 0.047 \\
SAE priority (low-income) & -0.058 & -0.088 & -0.103 & 0.019 & 0.026 & -0.008 & -0.064 \\
 & (0.003) & (0.003) & (0.004) & (0.002) & (0.004) & (0.004) & (0.004) \\
\bottomrule
\end{tabular}
}
\par\medskip\footnotesize
\begin{minipage}{\textwidth}
Notes: Same sample, scaling and controls as Table~\ref{tab:expected-va-sorting-v2} (comuna-by-cohort FE). "With priority" adds the SAE low-income (prioritario) and high-achievement flags. Because priority shifts admission probabilities for a given list, the change in the income coefficient in Panel B measures how much of the income gradient in offered VA runs through the priority rules.
\end{minipage}
\end{table}
